% Capitulo: A trilha hoje
\chapter{A trilha hoje}
\label{cap:trilha-hoje}

Quem começa hoje tem algo que nenhuma geração anterior teve: um tutor disponível a
qualquer hora, que explica um erro com paciência infinita, gera exercícios sob medida e
revisa uma consulta em segundos. É uma vantagem enorme, e também uma armadilha. O mesmo
assistente que explica pode simplesmente fazer, e quem deixa a IA fazer durante o
aprendizado termina o curso sem ter aprendido. Este capítulo trata das duas coisas: como
usar a IA para aprender mais rápido, e como não usá-la para deixar de aprender. No fim,
propõe uma ordem de estudo para quem está começando agora.

\section{A IA como tutor}
\label{sec:ia-tutor}

Os usos da IA que mais ajudam no aprendizado têm algo em comum: neles, quem está
aprendendo continua fazendo o trabalho de pensar. Alguns exemplos:

\begin{itemize}
  \item \textbf{Explicar um erro.} Em vez de pedir a consulta corrigida, colar a mensagem
    de erro e pedir uma explicação do que ela significa e de onde costuma vir.
  \item \textbf{Explicar código alheio.} Pedir que a IA explique, linha por linha, uma
    consulta ou um script herdado, e depois tentar reescrevê-lo sem olhar.
  \item \textbf{Gerar exercícios.} Pedir dez exercícios de junção, em dificuldade
    crescente, sobre um conjunto de dados conhecido, com as respostas guardadas para
    depois.
  \item \textbf{Revisar uma solução.} Escrever a consulta primeiro e só então pedir uma
    revisão: o que pode dar errado, que caso de borda ficou de fora, como ficaria mais
    legível.
  \item \textbf{Ser questionado.} Pedir que a IA faça perguntas sobre um assunto recém
    estudado, sem dar as respostas, e corrija o que estiver errado.
\end{itemize}

A evidência mais clara de que a forma de usar faz diferença vem de um experimento com
quase mil estudantes de ensino médio, em aulas de matemática. Um grupo usou um assistente
parecido com os de uso geral; outro usou um assistente configurado como tutor, que conhecia
a solução de cada exercício e os erros comuns, mas tinha instrução de dar dicas
progressivas sem entregar a resposta. Os dois grupos foram muito melhor que o grupo de
controle nos exercícios feitos com ajuda. Na prova, sem ajuda, o grupo do assistente comum
foi pior do que quem nunca teve acesso, e o grupo do tutor ficou no mesmo nível do controle
\parencite{bastani2025}. O assistente sem limites virou muleta; o tutor, não.

\begin{dica}
Ao estudar com IA, comece o pedido com algo como: \enquote{Estou aprendendo. Não me dê a
resposta. Me ajude a chegar nela com perguntas e dicas, uma de cada vez.} A diferença no
que fica aprendido é grande.
\end{dica}

\section{O risco do atalho}
\label{sec:risco-atalho}

O risco não é hipotético nem restrito a estudantes de ensino médio. Em 2026, um experimento
com programadores, na maioria em início de carreira, pediu que aprendessem uma biblioteca de
Python que não conheciam. Metade podia usar um assistente de IA; metade programou à mão. No
teste de compreensão feito depois, o grupo com IA acertou, em média, metade das questões;
o grupo sem IA, cerca de dois terços. A maior diferença apareceu justamente nas questões de
depuração, a habilidade necessária para perceber quando o código gerado está errado
\parencite{shen2026}. E o grupo com IA nem chegou a terminar a tarefa significativamente
mais rápido.

O mesmo estudo traz a parte animadora. Entre os que usaram IA, os que se saíram melhor
foram os que a usaram para entender, fazendo perguntas sobre o código, e não só para gerar
o código. A ferramenta era a mesma; o resultado dependeu de como foi usada.

Uma pesquisa com profissionais do conhecimento aponta na mesma direção, por outro caminho:
quanto mais as pessoas confiavam na IA para uma tarefa, menos esforço de pensamento crítico
diziam aplicar a ela; quanto mais confiavam na própria capacidade, mais esforço aplicavam
\parencite{lee2025}. O atalho é confortável, e o conforto é exatamente o problema.

A explicação para tudo isso é antiga. A pesquisa sobre aprendizagem mostra há décadas que
o esforço de recuperar uma informação, errar e corrigir é o que consolida o conhecimento
\parencite{bjork1994}, e que a habilidade de especialista vem de prática deliberada: tarefas
um pouco além da capacidade atual, com feedback, repetidas com atenção
\parencite{ericsson1993}. Quando a IA faz a parte difícil, ela remove justamente o esforço
que geraria o aprendizado. O resultado parece bom na hora e não fica.

\begin{destaque}
Usar IA para fazer é ótimo para quem já sabe. Para quem está aprendendo, é preciso usá-la
para entender, ou ela ensina só a pedir.
\end{destaque}

Isso não significa estudar sem IA. Significa separar dois momentos. No momento de
aprender uma coisa nova, a IA explica, pergunta e revisa, mas quem escreve é a pessoa. No
momento de produzir algo que já se sabe fazer, a IA pode escrever, e a pessoa revisa. O
erro é usar o segundo modo no primeiro momento.

\section{Uma ordem de estudo para hoje}
\label{sec:ordem-estudo}

Com essas ideias em mente, esta é a ordem que eu recomendaria para quem começa agora em
análise ou engenharia de dados. Ela não é muito diferente da trilha tradicional no
conteúdo. Muda na ordem, na profundidade de cada etapa e na forma de usar a IA em cada uma.

\begin{enumerate}
  \item \textbf{Dados em tabela e SQL, à mão.} Seleção, filtro, agregação, junções, funções
    de janela. É a base de tudo e a linguagem que mais tempo vai durar. A IA entra como
    tutor e gerador de exercícios, nunca como quem escreve as consultas.
  \item \textbf{Modelagem.} Grão, chaves, fatos e dimensões (capítulo~\ref{cap:dw}). Na
    trilha tradicional isso vinha tarde, quando vinha. Hoje vem cedo, porque é o que
    permite avaliar se o que a IA produziu faz sentido.
  \item \textbf{Python para dados.} Ler, limpar, transformar e visualizar dados com
    código. Aqui a IA já pode ajudar mais, desde que cada trecho gerado seja lido e
    entendido antes de ser usado.
  \item \textbf{Estatística aplicada e pensamento crítico.} Distribuições, comparação,
    incerteza, correlação e causalidade (\veja{sec:fund-estatistica}). É o que impede
    conclusões erradas tiradas de números certos.
  \item \textbf{Visualização e comunicação.} Escolher o gráfico, escrever a conclusão
    primeiro, apresentar para quem decide (\veja{sec:visualizacao}).
  \item \textbf{Práticas de engenharia.} Versionamento, testes, automação
    (Parte~\ref{parte:devops}). Antes eram território do engenheiro; hoje qualquer
    analista que trabalha com IA precisa delas para revisar e manter o que produz.
  \item \textbf{IA como ferramenta de produção.} Usar assistentes e agentes para
    acelerar o que já se sabe fazer, com critério para revisar e saber de que lado da
    fronteira está cada tarefa (\veja{sec:copiloto}).
\end{enumerate}

Em paralelo a tudo isso, desde o primeiro dia, aprender o negócio onde os dados nascem.
O capítulo~\ref{cap:negocio} explica por que essa é a parte que mais diferencia um
profissional, e o capítulo~\ref{cap:mapas} detalha a trilha para cada carreira.

\rascunho{
  Opcional: um parágrafo meu sobre como eu oriento hoje quem entra no time, ou o que
  observo de diferente entre quem aprendeu antes e quem está aprendendo com IA.
}

\begin{leituras}
  \leitura{bastani2025}{O experimento com estudantes que mostra a diferença entre usar a
    IA como muleta e como tutor. Leitura rápida e muito clara.}
  \leitura{shen2026}{Um experimento controlado sobre como a IA afeta a formação de
    habilidades de programação, com atenção especial à depuração.}
  \leitura{ericsson1993}{O artigo clássico sobre prática deliberada, a base de boa parte do
    que se sabe sobre como alguém se torna especialista.}
\end{leituras}
